\paragraph{Step 20 finite-carrier predicate statement.}
Let \(H\) be the declared finite carrier, \(q:H\to Q\) the current
access quotient, and \(s:H\to S\) a declared role readout with
\(\mathcal O_s\neq\emptyset\).  A candidate resolution record \(R\)
is evaluated by eight computable predicates
\[
\mathrm{Mem}(R),\mathrm{Hid}(R),\mathrm{Med}(R),\mathrm{Bud}(R),
\mathrm{Scp}(R),\mathrm{Crs}(R),\mathrm{Out}(R),\mathrm{Blk}(R).
\]
Each predicate is the conjunction of RoleSplit\((q,s)\), a unique
declared mechanism flag, and the corresponding operational witness:
memory residue record, hidden upstream latent, explicit admissible
mediation object, priced residual, proper scope restriction, target
coarsening, named outside-scope direction, or blocked nonclosure.

\paragraph{Exclusivity and coverage on the Step 20 suite.}
On the declared Step 20 finite carrier and candidate suite,
exactly one predicate fires for every candidate record, and every
family is selected by at least one candidate.  Thus the constructed
selectors are mutually exclusive and exhaustive on the suite used in
this diagnostic.

\paragraph{Classification of \(L\).}
For the candidate package \(L\), the mediation predicate fires because
the candidate record contains an explicit admissible object carrying
both \(q\) and \(s\), passing the native gates, commuting with the
declared projections, and strictly refining \(q\).  The other seven
predicates fail on their own operational criteria.  Therefore
\[
  L \mapsto \operatorname{BridgeMediatedRole}
\]
on the Step 20 finite carrier.
